% Generated draft: requires mathematical and source review.
\documentclass[11pt]{article}
\usepackage[margin=25mm]{geometry}
\usepackage{fontspec}
\setmainfont{Latin Modern Roman}
\usepackage{amsmath,amssymb,mathtools,unicode-math}
\setmathfont{Latin Modern Math}
\usepackage{enumitem,needspace,xurl,hyperref}
\setlength{\parindent}{0pt}
\setlength{\parskip}{5pt}
\setlength{\emergencystretch}{3em}
\newcommand{\N}{\mathbb{N}}
\newcommand{\Z}{\mathbb{Z}}
\newcommand{\Q}{\mathbb{Q}}
\newcommand{\R}{\mathbb{R}}
\newcommand{\C}{\mathbb{C}}
\newcommand{\F}{\mathbb{F}}
\newcommand{\A}{\mathbb{A}}
\newcommand{\PP}{\mathbb{P}}
\newcommand{\E}{\mathbb{E}}
\newcommand{\eps}{\varepsilon}
\newcommand{\dd}{\,\mathrm{d}}
\begin{document}
\section*{UnsolvedMath ID: 701046}
\textbf{Source problem code:} catalogue entry
\subsection{Definitions and mathematical statement}
Let \(\mathcal S=(P,B,I)\) be an incidence structure, where \(P\) is the set of points, \(B\) is the set of lines, and \(I\subseteq P\times B\) is the incidence relation. Write \(P_l=\{a\in P:(a,l)\in I\}\) and \(B_a=\{l\in B:(a,l)\in I\}\). Distinct points are collinear if they lie on a common line; put
\[
\Gamma(a)=\{b\in P\setminus\{a\}:B_a\cap B_b\ne\varnothing\}.
\]
A flag is an incident point--line pair, and an anti-flag is a nonincident point--line pair.

The supplied definition calls \(\mathcal S\) a semipartial geometry with parameters \((\alpha,s,t,\mu)\) when
\begin{align*}
|B_a|&=t+1 &&(a\in P),\\
|B_a\cap B_b|&\leq 1 &&(a,b\in P,\ a\ne b),\\
|P_l|&=s+1 &&(l\in B),\\
\bigl|\{m\in B_a:P_m\cap P_l\ne\varnothing\}\bigr|&\in\{0,\alpha\}
&&(a\in P,\ l\in B,\ (a,l)\notin I),\\
|\Gamma(a)\cap\Gamma(b)|&=\mu
&&(a,b\in P,\ a\ne b,\ b\notin\Gamma(a)).
\end{align*}
The source gives no explicit parameter ranges or finiteness, connectedness, or nondegeneracy assumptions; these are unresolved qualifications of the classification target.

An automorphism is a pair of permutations of \(P\) and \(B\) preserving incidence in both directions. A subgroup \(G\leq\operatorname{Aut}(\mathcal S)\) is flag-transitive if it acts transitively on the flags. In the usual algebraic meaning, \(G\) is almost simple if there is a nonabelian simple group \(T\) such that, identifying \(T\) with its inner automorphism group,
\[
T\leq G\leq\operatorname{Aut}(T).
\]
Whether the problem additionally requires \(T\), \(G\), and \(\mathcal S\) to be finite is not specified in the supplied record.

\textbf{Classification problem (Kourovka 14.61).} Determine all pairs \((\mathcal S,G)\) such that \(\mathcal S\) is a semipartial geometry as above and \(G\) is an almost simple flag-transitive group of automorphisms of \(\mathcal S\). The intended equivalence relation for classifying pairs is also unspecified.
\subsection{Short English statement}
Classify the pairs consisting of a semipartial geometry and an almost simple automorphism group acting transitively on its incident point--line pairs. The supplied incidence and common-neighbor conditions are retained, but finiteness, parameter restrictions, degeneracy conventions, and the equivalence relation for the classification require clarification.
\subsection{Sources}
\begin{itemize}
\item Local UnsolvedMath registry: supplied complete source record, permanent ID \(701046\), exact source problem code \texttt{catalogue entry}, titled \emph{Kourovka Notebook Problem 14.61}. The record supplies the status \texttt{open}; this is a registry label, not an independently verified mathematical or literature assessment. The record was created and updated on \texttt{2026-10-16T14:00:00Z}.
\item Attributed source: \emph{The Kourovka Notebook (Unsolved Problems in Group Theory)}, 14th issue (1999), Problem 14.61; proposed by A. A. Makhn\"{e}v in 1999. The supplied background also identifies the 21st issue edition (2026), arXiv identifier \texttt{1401.0300v48}. The supplied collection description names E. I. Khukhro and V. D. Mazurov as editors. These bibliographic details are supplied attribution and have not been independently checked.
\end{itemize}
\end{document}
